\subsection{Extension of positive linear forms}

PROPOSITION 1. --- Let E be a preordered vector space and V be a vector subspace of E such that every element of E is bounded above by an element of V. Given a linear form f on V that is positive for the preordered vector space structure of V (induced by that of E) there exists a non-empty set \(S_{f}\) of positive linear forms on E, each being an extension of f.~If \(h \in S_{f}\) then the values \(h(a)\) for \(a \in E\) lie in the interval \([\alpha', \alpha'']\) , where

\[
\alpha^ {\prime} = \sup _ {z \in \mathbf {V}, z \leqslant a} f (z), \quad \alpha^ {\prime \prime} = \inf _ {y \in \mathbf {V}, y \geqslant a} f (y).\tag{1}
\]

I. Special case.

Suppose firstly that \(E = V + R a\) . Since the proposition is trivial if \(a \in V\) , we confine ourselves to the case \(a \notin V\) . The conditions on V imply that the set \(A''\) of \(y \in V\) such that \(a \leqslant y\) is not empty; similarly the set \(A'\) of \(z \in V\) such that \(-z \geqslant -a\) (i.e.~\(z \leqslant a\) ) is not empty. For \(y \in A''\) and \(z \in A'\) , we have \(z \leqslant a \leqslant y\) , and thus by hypothesis \(f(z) \leqslant f(y)\) . Thus \(\alpha', \alpha''\) are finite and \(\alpha' \leqslant \alpha''\) . Any linear form \(f_1\) on E that extends f is completely determined by \(f_1(a)\) and for all \(\lambda \in R\) and all \(x \in V\) , we have

\[
f _ {1} (x + \lambda a) = f (x) + \lambda f _ {1} (a).
\]

Thus \(f_{1}\) is positive if and only if the relations

\[
x \in \mathbf {V}, \quad \lambda \in \mathbf {R}, \quad x + \lambda a \geqslant 0\tag{2}
\]

imply

\[
f (x) + \lambda f _ {1} (a) \geqslant 0.\tag{3}
\]

As \(f(\mu x) = \mu f(x)\) and the relations \(x \geqslant 0\) and \(\mu x \geqslant 0\) are equivalent for \(\mu > 0\) , it is sufficient to show that (2) implies (3) in the particular cases \(\lambda = 0, \lambda = 1\) and \(\lambda = -1\) . For \(\lambda = 0\) , the fact that (2) implies (3) follows from the hypothesis that \(f\) is positive. For \(\lambda = 1\) , to say that (2) implies (3) means that for \(-x \in A'\) , we have \(f_1(a) \geqslant f(-x)\) , i.e.~\(f_1(a) \geqslant \alpha'\) ; for \(\lambda = -1\) , (2) implies (3), means that for \(x \in A''\) , we have \(f(x) \geqslant f_1(a)\) , i.e.~\(f_1(a) \leqslant \alpha''\) . The proposition is therefore proved in this case.

\paragraph{II. General case.}

Let \(\mathfrak{F}\) be the set of pairs \((\mathrm{W}, g)\) where \(\mathrm{W}\) is a vector subspace of \(\mathrm{E}\) containing \(\mathrm{V}\) and \(g\) is a positive linear form on \(\mathrm{W}\) which is an extension of \(f\) . We order \(\mathfrak{F}\) putting \((\mathrm{W}, g) \leqslant (\mathrm{W}', g')\) if \(\mathrm{W} \subset \mathrm{W}'\) and if \(g'\) is an extension of \(g\) . Clearly \(\mathfrak{F}\) is inductive and by th. 2 of S, III, § 2.4, there is a maximal element \((\mathrm{W}_0, g_0)\) . Suppose \(\mathrm{W}_0 \neq \mathrm{E}\) . Then there exists a vector \(b \notin \mathrm{W}_0\) , and, if \(\mathrm{W}_1 = \mathrm{W}_0 + \mathbf{R}b\) , the special case above shows that there exists a positive linear form on \(\mathrm{W}_1\) which is an extension of \(g_0\) ; this contradicts the hypothesis that \((\mathrm{W}_0, g_0)\) is maximal. Thus \(\mathrm{W}_0 = \mathrm{E}\) , and the first part of the proposition is proved. When \(a \in \mathrm{V}\) , the second assertion is obviously true with \(\alpha' = \alpha'' = f(a)\) ; if, on the contrary, \(a \notin \mathrm{V}\) and one puts \(\mathrm{V}_1 = \mathrm{V} + \mathbf{R}a\) , the second assertion follows from the special case I of the proof.

COROLLARY. --- In a topological vector space E with a compatible preorder structure, let P be the set of elements \(\geqslant 0\) in E. Let V be a vector subspace of E containing at least one interior point \(x_{0}\) of P. Then every positive linear form on V can be extended to a positive linear form on E.

By prop. 1 it is sufficient to show that for every \(x \in \mathbf{E}\) , there exists \(x' \in \mathbf{V}\) such that \(x' - x \in \mathbf{P}\) . Now let \(\mathbf{U}\) be a neighbourhood of 0 in \(\mathbf{E}\) such that \(x_0 + \mathbf{U} \subset \mathbf{P}\) . Then \(x + x_0 + \mathbf{U} \subset x + \mathbf{P}\) , and, hence there exists \(\varepsilon\) such that \(0 < \varepsilon < 1\) and the point \(y = x_0 + (1 - \varepsilon)x\) belongs to \(x + \mathbf{P}\) ; then every point of the form \(x + \lambda(y - x)\) belongs to \(x + \mathbf{P}\) for \(\lambda > 0\) . If we take \(\lambda = 1 / \varepsilon\) , then \(x + \lambda(y - x) = \lambda x_0 \in \mathbf{V}\) , from which the conclusion follows.

The conclusion of the corollary is not necessarily valid if one does not assume that V contains an interior point of P, even if E is of finite dimension and if \(P \cap V\) contains points interior in V (II, p.~91, exerc. 25, b)).
% semantic-fragments: legacy-input-seam
\relax{}
